% sections/03-method.tex

\section{Method}
\label{sec:method}

We present three theorems forming a dependency chain:
Theorem~\ref{thm:sensitivity} derives the sensitivity bound;
Theorem~\ref{thm:mechanism} instantiates the private release mechanism using that bound;
Theorem~\ref{thm:composition} applies zCDP composition to the prospective setting.

% ── Theorem 1 ─────────────────────────────────────────────────────────────────

\subsection{Theorem 1: Sensitivity Bound}
\label{sec:method:sensitivity}

\begin{theorem}[Sensitivity Bound]
\label{thm:sensitivity}
Let $D$ and $D'$ be neighboring datasets differing in exactly one reporting node
inside window $w$ (i.e., one node's count changes by $\pm 1$).
The global sensitivity of the Kulldorff log-likelihood ratio satisfies:
\[
    \sens = \max_{D, D' : d(D,D')=1} |\LLR(w, D) - \LLR(w, D')| \leq f(\kw)
\]
where $f(\kw)$ is a closed-form function satisfying:
(i) $f$ is monotonically decreasing in $\kw$, and
(ii) $f(\kw) \to 0$ as $\kw \to \infty$.
\end{theorem}

\begin{proof}
\todo{Derive the closed-form upper bound $f(\kw)$.
Key step: when one node changes by $\pm 1$, show that $c$, $C$, and $\mu = C \cdot \nw / N$
all shift simultaneously.
Compute $|\LLR(c, C, \mu) - \LLR(c \pm 1, C \pm 1, \mu')|$ and find the maximum
over all valid $(c, C, \nw, \Ntot)$ with $\kw$ nodes.}
\end{proof}

\begin{remark}
The candidate form $f(\kw) = \ln(\kw / (\kw - 1)) + 1/\kw$ satisfies properties (i)
and (ii) above.
This form is \emph{not} the proven result; it is a heuristic starting point for the
derivation.
All numerical sensitivity values are withheld until Theorem~\ref{thm:sensitivity}
is formally proven.
\end{remark}

\begin{remark}
$\kw$ does not appear in the LLR formula~\eqref{eq:llr} because it is irrelevant to
outbreak detection.
It is essential to privacy: it determines how many independent actors can each
change $\cobs$ by $\pm 1$.
The simultaneous change in $\Ctot$ and $\muexp$ (which depends on $\Ctot$) is what
makes the derivation non-trivial and propagates globally to all candidate windows.
\end{remark}

% ── Theorem 2 ─────────────────────────────────────────────────────────────────

\subsection{Theorem 2: Private Window Selection}
\label{sec:method:mechanism}

\begin{theorem}[Private Window Selection]
\label{thm:mechanism}
Define the normalized utility function:
\[
    \util = \frac{\LLR(w)}{\sens}
\]
The global sensitivity of $u$ is $1$.
The mechanism that selects window $\wstar$ with probability:
\[
    \Pr[\wstar = w] \propto \exp\!\left(\frac{\eps \cdot \util}{2}\right)
\]
is $\eps$-differentially private.

Furthermore, let $w_0 = \arg\max_w \util$ be the \emph{privacy-adjusted optimal window}
and define:
\[
    \gapu = u(w_0) - \max_{w \neq w_0} u(w) = \frac{\LLR(w_0)}{\Delta(k_{w_0})} - \max_{w \neq w_0}\frac{\LLR(w)}{\Delta(k_w)}
\]
Then:
\[
    \Pr[\wstar = w_0] \geq 1 - |\Wcal| \cdot \exp\!\left(\frac{-\eps \cdot \gapu}{2}\right)
\]
\end{theorem}

\begin{proof}
\textit{(Global sensitivity = 1.)}
When one node's count changes by $\pm 1$, $\LLR(w)$ changes by at most $\sens$
(Theorem~\ref{thm:sensitivity}).
Therefore $u(w) = \LLR(w)/\sens$ changes by at most $\sens/\sens = 1$.

\textit{($\eps$-DP and utility guarantee.)}
The mechanism is the exponential mechanism~\cite{mcsherry2007} instantiated with
utility function $u$ and global sensitivity $\Delta u = 1$.
$\eps$-DP follows directly from the exponential mechanism privacy guarantee.
The utility bound follows from Theorem~\ref{thm:mst-utility} with $\Delta q = 1$
and $\Delta = \gapu$.
\end{proof}

\begin{remark}
$w_0$ is the \emph{privacy-adjusted optimal window}: the candidate with the highest
ratio of statistical evidence to sensitivity.
This differs from both the ground-truth outbreak location and the non-private argmax
(highest raw LLR).
A dense window with moderate LLR may outrank a sparse window with high LLR because
$\sens$ is smaller for dense windows.
\end{remark}

\begin{remark}
The Monte Carlo p-value (Step~3 of Algorithm~\ref{alg:kulldorff}) is applied to $\wstar$
after the DP mechanism runs.
By the post-processing theorem~\cite[Proposition~2.1]{dwork2013}, this costs no
additional privacy budget.
\end{remark}

Algorithm~\ref{alg:private} presents the complete private scan statistic procedure.
It modifies only Step~2 of Algorithm~\ref{alg:kulldorff}; Steps~1 and~3 are unchanged.

% algorithms/alg2-private.tex
% Algorithm 2: Private scan statistic — YOUR contribution (Theorem 2)
% Placed in the method section. Differs from Algorithm 1 only in Step 2.

\begin{algorithm}[H]
\caption{Private Prospective Scan Statistic (Theorem~\ref{thm:mechanism})}
\label{alg:private}
\begin{algorithmic}[1]

\Statex \textbf{Input:}
  Same as Algorithm~\ref{alg:kulldorff}, plus privacy budget $\eps$
\Statex \textbf{Output:}
  Privately selected cluster $\wstar$ with $\LLR(\wstar)$, $\RR(\wstar)$, $p$-value

\Statex
\Statex \textit{// Step 1: Score every candidate window (identical to Algorithm~\ref{alg:kulldorff})}
\State Run Step~1 of Algorithm~\ref{alg:kulldorff} to obtain $\LLR(w)$ for all $w \in \Wcal$

\Statex
\Statex \textit{// Step 2: Private window selection (replaces deterministic argmax)}
\For{each $w \in \Wcal$}
    \State $\kw \gets$ number of reporting nodes inside $w$
    \State $\sens \gets f(\kw)$
        \Comment{sensitivity bound from Theorem~\ref{thm:sensitivity}}
    \State $u(w) \gets \LLR(w) / \sens$
        \Comment{normalized utility, global sensitivity $= 1$}
    \State $p_w \gets \exp\!\left(\eps \cdot u(w) / 2\right)$
        \Comment{unnormalized selection weight}
\EndFor
\State $Z \gets \sum_{w \in \Wcal} p_w$
    \Comment{partition function}
\State Sample $\wstar \sim \Pr[\wstar = w] = p_w / Z$
    \Comment{exponential mechanism~\cite{mcsherry2007}}

\Statex
\Statex \textit{// Step 3: Significance testing (post-processing; no additional budget)}
\State Run Step~3 of Algorithm~\ref{alg:kulldorff} on $\wstar$
    \Comment{safe by post-processing theorem~\cite[Prop.~2.1]{dwork2013}}

\State \Return $\wstar,\; \LLR(\wstar),\; \RR(\wstar),\; p$

\end{algorithmic}
\end{algorithm}


% ── Theorem 3 ─────────────────────────────────────────────────────────────────

\subsection{Theorem 3: Prospective Composition}
\label{sec:method:composition}

In the prospective surveillance setting~\cite{kulldorff2001}, the scan statistic is
run once per week for $T$ consecutive weeks, releasing a cluster alert each time.

\begin{theorem}[Prospective Composition]
\label{thm:composition}
Let the mechanism from Theorem~\ref{thm:mechanism} be run for $T$ weeks, each release
contributing $\rhow$-zCDP.
Then:
\[
    \rhotot = T \cdot \rhow
\]
Converting to $(\eps, \delta)$-DP via~\cite{bun2016}:
\[
    \eps_{\mathrm{total}} = T\rhow + 2\sqrt{T\rhow \cdot \ln(1/\delta)}
\]
For small $\rhow$, the maximum number of weeks within a target budget $\eps_{\mathrm{target}}$
is approximately:
\[
    \Tmax \approx \frac{\eps_{\mathrm{target}}^2}{4\rhow \cdot \ln(1/\delta)}
\]
\end{theorem}

\begin{proof}
\todo{Fill in the proof. Use linear composition of $\rho$-zCDP mechanisms
(\cite{bun2016} Lem 1.7, p.7; adaptive form Lem 2.3, p.12) to obtain
$\rho_{\mathrm{total}} = T\rhow$, then apply the zCDP~$\to$~$(\eps,\delta)$-DP
conversion (\cite{bun2016} Prop 1.3, p.6). Each weekly release is a fresh
invocation of the per-week $\rho$-zCDP mechanism on that week's dataset; zCDP
composition is over mechanisms, not over disjoint inputs, so temporal window
overlap does not affect the bound. For $T_{\max}$, drop the linear term $T\rhow$
(valid when $\rhow \ll \ln(1/\delta)$) and solve
$\eps \approx 2\sqrt{T\rhow\ln(1/\delta)}$ for $T$.}
\end{proof}

\begin{remark}
Under naive composition, $T$ releases each costing $\eps_{\mathrm{week}}$-DP accumulate
$T \cdot \eps_{\mathrm{week}}$ total.
With $\eps_{\mathrm{week}} = 0.1$ and a system budget of $\eps = 1$, this exhausts in
$T = 10$ weeks.
Under Theorem~\ref{thm:composition} with $\rhow = 0.01$, $\delta = 10^{-6}$:
$\Tmax \approx 1 / (4 \times 0.01 \times 13.8) \approx 180$ weeks ($\approx 3.5$ years).
zCDP composition is what makes continuous prospective surveillance deployable.
\end{remark}
